% Catalogo de figuras propuestas para el capitulo 2 de la tesis (version 2).
% Compilar: pdflatex propuesta-figuras.tex (dos pasadas). No va en la tesis.
\documentclass[a4paper,10pt]{article}
\usepackage[utf8]{inputenc}
\usepackage[spanish,es-tabla]{babel}
\usepackage[margin=2.2cm]{geometry}
\usepackage{amsmath,amssymb}
\usepackage{xcolor}
\usepackage{tikz}
\usetikzlibrary{arrows.meta,positioning,fit,backgrounds,calc,shapes.geometric}
\usepackage{pgfplots}
\pgfplotsset{compat=1.18}
\usepackage{subcaption}
\usepackage{caption}
\captionsetup{font=small,labelfont=bf}
\usepackage{float}
\floatplacement{figure}{H}

\definecolor{azul}{HTML}{4C78A8}
\definecolor{naranja}{HTML}{F58518}
\definecolor{verde}{HTML}{54A24B}
\definecolor{rojo}{HTML}{E45756}
\definecolor{violeta}{HTML}{B279A2}
\definecolor{gris}{HTML}{8C8C8C}

\newcommand{\donde}[1]{\par\noindent\textbf{Dónde va.} #1\par\smallskip}
\newcommand{\pie}[1]{\par\noindent\textbf{Pie propuesto.} #1\par\smallskip}
\newcommand{\nota}[1]{\par\noindent\textbf{Nota.} #1\par\smallskip}

\title{Figuras propuestas para el Capítulo 2 (versión 3)}
\author{Catálogo para evaluar, no forma parte de la tesis}
\date{Octubre de 2026}

\begin{document}
\maketitle

\noindent
Tercera vuelta. Quedaron afuera ``quién se mueve'', el perfil con datos, el
mapa del capítulo y el ejemplo del término cuadrático. Se rehicieron el
programa lineal, la planta, las zonas ABC y las dos fases, y se corrigieron el
costo de acceso y la leyenda del perfil esquemático. Paleta del manual, azul
\tikz\fill[azul] (0,0) rectangle (0.8em,0.8em); y naranja
\tikz\fill[naranja] (0,0) rectangle (0.8em,0.8em); sobre gris.

\tableofcontents

% =====================================================================
\section{Nuevas}
% =====================================================================

\subsection{N1. Programa lineal y relajación entera}

\donde{Nueva, en 2.7 Formulación matemática, después del párrafo que presenta
la programación lineal y el símplex. El párrafo la referenciaría en ``hay al
menos un vértice que lo alcanza'' y en ``los vértices del poliedro relajado ya
no tienen por qué ser enteros''.}

\begin{figure}
    \centering
    \begin{subfigure}{0.48\textwidth}
        \centering
        \begin{tikzpicture}[x=1.1cm, y=1.1cm]
            % region factible: x>=0, y>=0, x+2y<=7, 3x+y<=9
            \fill[azul!30] (0,0) -- (3,0) -- (2.2,2.4) -- (0,3.5) -- cycle;
            \draw[azul!80!black, very thick] (0,0) -- (3,0) -- (2.2,2.4) -- (0,3.5) -- cycle;
            % restricciones rotuladas sobre su arista
            \path (0,3.5) -- (2.2,2.4) node[pos=0.45, above, sloped, font=\scriptsize, black!60] {$x + 2y \le 7$};
            \path (3,0) -- (2.2,2.4) node[pos=0.5, above, sloped, font=\scriptsize, black!60] {$3x + y \le 9$};
            % ejes con marcas
            \draw[-{Stealth}, black!70] (-0.3,0) -- (4.2,0) node[below, font=\scriptsize] {$x$};
            \draw[-{Stealth}, black!70] (0,-0.3) -- (0,4.1) node[left, font=\scriptsize] {$y$};
            \foreach \t in {1,2,3} { \draw[black!70] (\t,0.06) -- (\t,-0.06) node[below, font=\scriptsize] {\t}; \draw[black!70] (0.06,\t) -- (-0.06,\t) node[left, font=\scriptsize] {\t}; }
            % curvas de nivel del objetivo x+3y = c
            \begin{scope}
                \clip (-0.2,-0.2) rectangle (4.1,4.0);
                \draw[naranja, dashed] (0,1.5) -- (4.1,0.1333);
                \draw[naranja, dashed] (0,2.5) -- (4.1,1.1333);
                \draw[naranja, dashed] (0,3.5) -- (4.1,2.1333);
            \end{scope}
            \node[font=\scriptsize, naranja, anchor=west] at (3.0,2.75) {$x + 3y = 10{,}5$};
            \node[font=\scriptsize, naranja, anchor=south west] at (3.15,1.5) {$x + 3y = 7{,}5$};
            \node[font=\scriptsize, naranja, anchor=south west] at (3.15,0.5) {$x + 3y = 4{,}5$};
            % direccion de mejora
            \draw[-{Stealth}, naranja, very thick] (1.1,0.5) -- ++(0.3,0.9) node[right, font=\scriptsize, pos=0.75] {mejora};
            % vertices
            \foreach \p in {(0,0),(3,0),(2.2,2.4)} \fill[azul!80!black] \p circle (0.07);
            \fill[naranja] (0,3.5) circle (0.1);
            \node[font=\scriptsize, anchor=west] at (0.18,3.72) {óptimo $(0;\,3{,}5)$};
        \end{tikzpicture}
        \caption{El programa lineal}
    \end{subfigure}
    \hfill
    \begin{subfigure}{0.48\textwidth}
        \centering
        \begin{tikzpicture}[x=1.1cm, y=1.1cm]
            \fill[azul!30] (0,0) -- (3,0) -- (2.2,2.4) -- (0,3.5) -- cycle;
            \draw[azul!80!black, very thick] (0,0) -- (3,0) -- (2.2,2.4) -- (0,3.5) -- cycle;
            \draw[-{Stealth}, black!70] (-0.3,0) -- (4.2,0) node[below, font=\scriptsize] {$x$};
            \draw[-{Stealth}, black!70] (0,-0.3) -- (0,4.1) node[left, font=\scriptsize] {$y$};
            \foreach \t in {1,2,3} { \draw[black!70] (\t,0.06) -- (\t,-0.06) node[below, font=\scriptsize] {\t}; \draw[black!70] (0.06,\t) -- (-0.06,\t) node[left, font=\scriptsize] {\t}; }
            % puntos enteros factibles, bien visibles
            \foreach \x in {0,...,3} \foreach \y in {0,...,3} {
                \pgfmathparse{(\x+2*\y<=7) && (3*\x+\y<=9) ? 1 : 0}
                \ifnum\pgfmathresult=1 \fill[azul!80!black] (\x,\y) circle (0.085); \fi
            }
            % nivel por el optimo relajado y por el mejor entero
            \begin{scope}
                \clip (-0.2,-0.2) rectangle (4.1,4.0);
                \draw[naranja, dashed] (0,3.5) -- (4.1,2.1333);
                \draw[naranja, dashed, thin] (0,3.3333) -- (4.1,1.9667);
            \end{scope}
            \draw[naranja, very thick] (0,3.5) circle (0.12);
            \node[font=\scriptsize, anchor=west, naranja] at (0.2,3.78) {óptimo relajado, $10{,}5$};
            \fill[naranja] (1,3) circle (0.12);
            \node[font=\scriptsize, anchor=west] at (1.2,3.36) {mejor entero $(1;\,3)$, $10$};
            \draw[black!70, very thick] (0,3) circle (0.12);
            \node[font=\scriptsize, anchor=west, black!70] at (0.2,2.55) {redondeo, $9$};
        \end{tikzpicture}
        \caption{Con variables enteras}
    \end{subfigure}
\end{figure}

\pie{Un programa lineal en dos variables y su versión entera. En (a) la región
factible es el polígono que recortan las restricciones, las líneas punteadas
son curvas de nivel del objetivo $x + 3y$, la flecha indica hacia dónde mejora
y el máximo se alcanza en un vértice. En (b) solo valen los puntos de
coordenadas enteras, marcados en azul. El vértice óptimo de la relajación no es
entero, el mejor punto entero no es un vértice del polígono, y tampoco es el
redondeo del óptimo relajado, que vale menos.}

\clearpage
\subsection{N2. Planta del depósito con el camino de un producto}

\donde{Reemplaza la Figura 2.2 actual, en 2.3.1, después del párrafo que define
las seis etapas. Ahora es un solo camino continuo, el de un producto, numerado
de 1 a 6.}

\begin{figure}
    \centering
    \begin{tikzpicture}[x=1cm, y=1cm, font=\scriptsize,
        num/.style={circle, fill=white, draw=black!70, inner sep=1pt, minimum size=3.8mm, font=\scriptsize}]
        % edificio
        \draw[black!70, thick] (0,0) rectangle (11,6.4);
        % puertas del dock
        \foreach \x in {1.2,2.2,3.2} \draw[fill=naranja!30, draw=naranja!70!black] (\x-0.35,-0.02) rectangle ++(0.7,0.18);
        \foreach \x in {7.8,8.8,9.8} \draw[fill=azul!30, draw=azul!70!black] (\x-0.35,-0.02) rectangle ++(0.7,0.18);
        \node[anchor=north] at (2.2,-0.12) {\emph{dock} de recepción};
        \node[anchor=north] at (8.8,-0.12) {\emph{dock} de despacho};
        % zonas
        \draw[fill=naranja!10, draw=naranja!60!black, dashed] (0.4,0.4) rectangle (4.0,1.5);
        \node at (2.2,0.95) {recepción};
        \draw[fill=azul!10, draw=azul!60!black, dashed] (7.0,0.4) rectangle (10.6,1.5);
        \node at (8.8,0.95) {empaque};
        % estanterias
        \foreach \i in {0,...,6} \draw[fill=gray!25, draw=black!60] (1.0+\i*1.5, 2.3) rectangle ++(0.5, 3.4);
        \node[anchor=south] at (9.5,5.8) {estanterías};
        % ubicacion donde queda el producto (cara de la estanteria 3, pasillo de x=3.5)
        \fill[naranja] (3.92,3.85) rectangle (4.1,4.15);
        % 1 -> 2: recepcion y put-away (naranja), sube por el lado derecho del pasillo
        \draw[-{Stealth}, very thick, naranja] (2.2,0.18) -- (2.2,0.4);
        \draw[-{Stealth}, very thick, naranja] (2.2,1.5) -- (2.2,2.05) -- (3.65,2.05) -- (3.65,4.0) -- (3.88,4.0);
        % 4: recorrido de picking (azul): sale de empaque, va por el corredor frontal,
        % sube por el lado izquierdo del mismo pasillo, cruza por el fondo y baja por
        % otro pasillo directo a empaque. Sin tramos repetidos.
        \fill[azul] (7.3,1.5) circle (0.07);
        \draw[-{Stealth}, very thick, azul] (7.3,1.5) -- (7.3,1.75) -- (3.35,1.75) -- (3.35,6.0) -- (8.0,6.0) -- (8.0,1.5);
        % 6: despacho
        \draw[-{Stealth}, very thick, azul] (8.8,0.4) -- (8.8,0.18);
        % numeros
        \node[num] at (0.75,-0.3) {1};
        \node[num] at (1.95,2.3) {2};
        \node[num] at (3.78,4.4) {3};
        \node[num] at (3.35,4.9) {4};
        \node[num] at (10.3,1.25) {5};
        \node[num] at (10.25,-0.3) {6};
        % referencias
        \draw[very thick, naranja] (0.0,6.75) -- (0.5,6.75) node[anchor=west, black] {recepción y \emph{put-away}};
        \draw[very thick, azul] (3.6,6.75) -- (4.1,6.75) node[anchor=west, black] {\emph{picking}, empaque y despacho};
    \end{tikzpicture}
\end{figure}

\pie{Planta esquemática de un depósito y el camino de un producto. Entra por el
\emph{dock} de recepción (1), el \emph{put-away} lo lleva hasta una ubicación
de las estanterías (2), donde queda almacenado (3) hasta que un recorrido de
\emph{picking}, que sale de empaque (punto azul), pasa a buscarlo (4) y lo deja
en empaque (5), de donde sale por el \emph{dock} de despacho (6). En naranja el
tramo de entrada, en azul el de salida.}

\nota{Flujo pasante, entrada de un lado y salida del otro, para que las etapas
se lean de izquierda a derecha. El recorrido de \emph{picking} sale de empaque
(punto azul), va por el corredor frontal, sube por el pasillo donde está el
producto, cruza por el fondo y baja por otro pasillo directo a empaque, que es
lo que hace la política en $S$ de 2.3.3. El \emph{put-away} sube por el mismo
pasillo, por eso las dos líneas van paralelas en ese tramo.}

\clearpage
\subsection{N3. Conveniencia de las ubicaciones y zonas ABC}

\donde{Nueva, en 2.5.4 Clases y restricciones, junto a la oración ``Las zonas
pueden dibujarse de dos maneras''. El tono es continuo, como en el manual, y
las clases se cortan sobre ese orden.}

\begin{figure}
    \centering
    \begin{subfigure}{0.46\textwidth}
        \centering
        \begin{tikzpicture}[x=1cm, y=1cm, font=\scriptsize]
            \path (-0.5,-0.7) rectangle (4.9,3.9);
            \foreach \i in {0,...,5} {
                \shade[top color=azul!8, bottom color=azul!90] (\i*0.7, 0.4) rectangle ++(0.3, 2.8);
                \draw[black!50] (\i*0.7, 0.4) rectangle ++(0.3, 2.8);
            }
            \draw[black!60, dashed] (-0.2,1.3) -- (4.0,1.3);
            \draw[black!60, dashed] (-0.2,2.2) -- (4.0,2.2);
            \node[anchor=west] at (4.0,0.85) {A};
            \node[anchor=west] at (4.0,1.75) {B};
            \node[anchor=west] at (4.0,2.7) {C};
            \fill[naranja] (1.9, 0.05) circle (0.08);
            \node[anchor=north] at (1.9,-0.05) {\emph{dock}};
        \end{tikzpicture}
        \caption{Por franjas, a lo largo de cada pasillo}
    \end{subfigure}
    \hfill
    \begin{subfigure}{0.46\textwidth}
        \centering
        \begin{tikzpicture}[x=1cm, y=1cm, font=\scriptsize]
            \path (-0.5,-0.7) rectangle (4.9,3.9);
            % un tono plano por pasillo: mas oscuro cuanto mas cerca del dock
            \foreach \i/\c in {0/azul!15, 1/azul!50, 2/azul!90, 3/azul!90, 4/azul!50, 5/azul!15} {
                \fill[\c] (\i*0.7, 0.4) rectangle ++(0.3, 2.8);
                \draw[black!50] (\i*0.7, 0.4) rectangle ++(0.3, 2.8);
            }
            \draw[black!60, dashed] (0.5,0.2) -- (0.5,3.4);
            \draw[black!60, dashed] (1.2,0.2) -- (1.2,3.4);
            \draw[black!60, dashed] (2.6,0.2) -- (2.6,3.4);
            \draw[black!60, dashed] (3.3,0.2) -- (3.3,3.4);
            \node[anchor=south] at (1.9,3.4) {A};
            \node[anchor=south] at (0.85,3.4) {B};
            \node[anchor=south] at (2.95,3.4) {B};
            \node[anchor=south] at (0.15,3.4) {C};
            \node[anchor=south] at (3.65,3.4) {C};
            \fill[naranja] (1.9, 0.05) circle (0.08);
            \node[anchor=north] at (1.9,-0.05) {\emph{dock}};
        \end{tikzpicture}
        \caption{Por pasillos enteros}
    \end{subfigure}

    \vspace{0.6em}
    \begin{tikzpicture}[font=\scriptsize]
        \shade[left color=azul!90, right color=azul!8] (0,0) rectangle (3,0.25);
        \draw[black!50] (0,0) rectangle (3,0.25);
        \node[anchor=east] at (-0.1,0.125) {más conveniente};
        \node[anchor=west] at (3.1,0.125) {menos conveniente};
    \end{tikzpicture}
\end{figure}

\pie{Dos maneras de cortar en clases las ubicaciones de un mismo depósito, con
el \emph{dock} abajo. El tono indica la conveniencia, más oscuro cuanto más
corto es el viaje de ida y vuelta desde el \emph{dock}, y las líneas punteadas
son los cortes entre las clases A, B y C. En (a) la conveniencia crece hacia el
frente de cada pasillo y las clases son franjas. En (b) lo que cuenta es el
pasillo, los centrales son más convenientes que los de los costados, y las
clases son pasillos enteros.}

\clearpage
\subsection{N4. Agrupar primero, asignar después}

\donde{Nueva, en 2.8 Estado del arte, en el bloque ``Método de asignación'',
donde se dice que la norma es la descomposición en dos fases.}

\begin{figure}
    \centering
    \begin{tikzpicture}[x=1cm, y=1cm, font=\scriptsize,
        chip/.style={draw=black!50, fill=gray!15, rounded corners=1pt, inner sep=2pt, minimum width=4.2mm, minimum height=4.2mm, font=\scriptsize},
        chipA/.style={chip, fill=azul!25, draw=azul!70!black},
        chipB/.style={chip, fill=naranja!30, draw=naranja!70!black},
        grupo/.style={draw, rounded corners=3pt, inner sep=4pt},
        panel/.style={font=\scriptsize\bfseries, anchor=south}]
        % --- panel 1: pedidos
        \node[panel] at (0.9,3.0) {pedidos};
        \foreach \y/\lista in {2.4/{A,B}, 1.8/{A,B,E}, 1.2/{C,D}, 0.6/{C,D}, 0.0/{F}} {
            \node[anchor=east] at (-0.1,\y) {};
            \foreach [count=\k] \s in \lista { \node[chip] at (0.05+\k*0.55,\y) {$\s$}; }
        }
        \node[anchor=north, black!60] at (0.9,-0.35) {$\cdots$};
        % --- flecha 1
        \draw[-{Stealth}, thick, black!60] (2.6,1.2) -- (3.7,1.2) node[midway, above, align=center] {medir afinidad\\y agrupar};
        % --- panel 2: grupos
        \node[panel] at (5.3,3.0) {grupos afines};
        \node[chipA] (a) at (4.4,2.2) {$A$};
        \node[chipA] (b) at (5.0,2.2) {$B$};
        \node[chipA] (e) at (5.6,2.2) {$E$};
        \node[grupo, draw=azul!70!black, fit=(a)(b)(e)] (g1) {};
        \node[chipB] (c) at (4.7,1.1) {$C$};
        \node[chipB] (d) at (5.3,1.1) {$D$};
        \node[grupo, draw=naranja!70!black, fit=(c)(d)] (g2) {};
        \node[chip] (f) at (5.0,0.1) {$F$};
        \node[grupo, draw=black!50, fit=(f)] (g3) {};
        % --- flecha 2
        \draw[-{Stealth}, thick, black!60] (6.9,1.2) -- (7.9,1.2) node[midway, above] {asignar};
        % --- panel 3: deposito
        \node[panel] at (9.8,3.0) {depósito};
        \foreach \i in {0,...,3} \draw[fill=gray!20, draw=black!60] (8.5+\i*0.8, 0.0) rectangle ++(0.35, 2.6);
        \fill[azul!60] (8.5,0.0) rectangle (8.85,1.1);
        \fill[naranja!70] (9.3,0.0) rectangle (9.65,0.8);
        \fill[gray!60] (10.9,2.1) rectangle (11.25,2.6);
        \fill[naranja] (9.1,-0.3) circle (0.07);
        \node[anchor=north] at (9.1,-0.38) {\emph{dock}};
        % referencias de color
        \fill[azul!60] (8.3,-1.05) rectangle ++(0.3,0.22); \node[anchor=west] at (8.65,-0.94) {$A$, $B$, $E$};
        \fill[naranja!70] (10.1,-1.05) rectangle ++(0.3,0.22); \node[anchor=west] at (10.45,-0.94) {$C$, $D$};
        \fill[gray!60] (11.5,-1.05) rectangle ++(0.3,0.22); \node[anchor=west] at (11.85,-0.94) {$F$};
    \end{tikzpicture}
\end{figure}

\pie{El esquema de dos fases que domina la literatura. Del historial de pedidos
se mide qué SKU aparecen juntas y con eso se arman grupos de SKU afines. Recién
después se decide dónde va cada grupo en el depósito, los grupos con más
actividad cerca del \emph{dock} y las SKU sin pareja donde quede lugar. Las dos
decisiones se toman por separado, y lo que se pierde al separarlas rara vez se
mide.}

% =====================================================================
\clearpage
\section{Existentes, recoloreadas}
% =====================================================================

\noindent
R1 a R4 quedaron como en la primera versión, que te gustaron. R5 cambia la
forma de dibujar los viajes y R6 arregla la leyenda.

\subsection{R1. Figura 2.1, la unidad de manipulación}
\donde{2.1 Evolución, misma posición.}
\begin{figure}
    \centering
    \begin{tikzpicture}[x=1cm, y=1cm, every node/.style={font=\small}]
        \foreach \i in {0,1,2} \foreach \j in {0,1,2}
            \draw[fill=azul!18, draw=azul!70!black] (\i*0.63, \j*0.5) rectangle ++(0.58, 0.44);
        \draw[fill=gray!55, draw=black!70] (-0.08, -0.2) rectangle (1.92, -0.04);
        \node at (0.92, -0.65) {\emph{pallet}};
        \draw[-{Stealth}, thick, black!60] (2.25, 0.7) -- (2.95, 0.7);
        \foreach \i in {0,1} \foreach \j in {0,1}
            \draw[fill=azul!18, draw=azul!70!black] (3.25+\i*0.46, 0.24+\j*0.46) rectangle ++(0.42, 0.42);
        \draw[draw=black!70, thick] (3.2, 0.19) rectangle (4.22, 1.21);
        \node at (3.71, -0.65) {caja};
        \draw[-{Stealth}, thick, black!60] (4.55, 0.7) -- (5.25, 0.7);
        \foreach \i in {0,1,2} \draw[fill=azul!18, draw=azul!70!black] (5.55+\i*0.36, 0.54) rectangle ++(0.32, 0.32);
        \draw[draw=black!70, thick] (5.5, 0.49) rectangle (6.66, 0.91);
        \node at (6.08, -0.65) {paquete interno};
        \draw[-{Stealth}, thick, black!60] (7.0, 0.7) -- (7.7, 0.7);
        \draw[fill=azul!18, draw=azul!70!black] (7.95, 0.54) rectangle ++(0.32, 0.32);
        \node at (8.11, -0.65) {unidad};
    \end{tikzpicture}
\end{figure}

\subsection{R2. Figura 2.3, variantes de geometría}
\donde{2.3.2 La geometría del depósito, misma posición.}
\begin{figure}
    \centering
    \begin{subfigure}{0.48\textwidth}
        \centering
        \begin{tikzpicture}[x=1cm, y=1cm]
            \foreach \i in {0,...,9} \draw[fill=gray!15, draw=black!60] (\i*0.42, 0) rectangle ++(0.18, 2.6);
            \fill[azul, opacity=0.18] (1.10, 0) rectangle (2.86, 2.6);
            \fill[azul, opacity=0.22] (1.52, 0) rectangle (2.44, 2.6);
            \fill[naranja] (1.98, -0.22) circle (0.07);
            \fill[naranja] (1.98, 2.82) circle (0.07);
            \node[font=\scriptsize] at (1.98, -0.5) {recepción};
            \node[font=\scriptsize] at (1.98, 3.08) {despacho};
        \end{tikzpicture}
        \caption{Flujo pasante}
    \end{subfigure}
    \hfill
    \begin{subfigure}{0.48\textwidth}
        \centering
        \begin{tikzpicture}[x=1cm, y=1cm]
            \foreach \i in {0,...,9} \draw[fill=gray!15, draw=black!60] (\i*0.42, 0) rectangle ++(0.18, 2.6);
            \fill[azul, opacity=0.14] (1.98, 2.35) -- (0.15, 0) -- (3.81, 0) -- cycle;
            \fill[azul, opacity=0.18] (1.98, 1.5) -- (0.85, 0) -- (3.11, 0) -- cycle;
            \fill[azul, opacity=0.24] (1.98, 0.75) -- (1.45, 0) -- (2.51, 0) -- cycle;
            \fill[naranja] (1.98, -0.22) circle (0.07);
            \node[font=\scriptsize] at (1.98, -0.5) {recepción y despacho};
            \node[font=\scriptsize, opacity=0] at (1.98, 3.08) {despacho};
        \end{tikzpicture}
        \caption{Flujo en U}
    \end{subfigure}

    \vspace{0.8em}

    \begin{subfigure}{0.48\textwidth}
        \centering
        \begin{tikzpicture}[x=1cm, y=1cm]
            \foreach \i in {0,...,9} {
                \draw[fill=gray!15, draw=black!60] (\i*0.42, 1.5) rectangle ++(0.18, 1.1);
                \draw[fill=gray!15, draw=black!60] (\i*0.42, 0) rectangle ++(0.18, 1.1);
            }
            \draw[very thick, azul] (0.72, 0.3) -- (0.72, 1.3) -- (3.24, 1.3) -- (3.24, 2.25);
            \fill[naranja] (0.63, 0.16) rectangle ++(0.18, 0.18);
            \fill[naranja] (3.15, 2.2) rectangle ++(0.18, 0.18);
            \fill[naranja] (1.98, -0.22) circle (0.07);
            \node[font=\scriptsize] at (1.98, -0.5) {dock};
        \end{tikzpicture}
        \caption{Pasillo transversal}
    \end{subfigure}
    \hfill
    \begin{subfigure}{0.48\textwidth}
        \centering
        \begin{tikzpicture}[x=1cm, y=1cm]
            \begin{scope}
                \clip (0, 0) -- (1.68, 0) -- (0.08, 2.6) -- (0, 2.6) -- cycle;
                \foreach \j in {0,...,5} \draw[fill=gray!15, draw=black!60] (0, 0.06+\j*0.44) rectangle (1.9, 0.24+\j*0.44);
            \end{scope}
            \begin{scope}
                \clip (2.28, 0) -- (3.96, 0) -- (3.96, 2.6) -- (3.88, 2.6) -- cycle;
                \foreach \j in {0,...,5} \draw[fill=gray!15, draw=black!60] (2.0, 0.06+\j*0.44) rectangle (3.96, 0.24+\j*0.44);
            \end{scope}
            \begin{scope}
                \clip (1.98, 0.2) -- (0.5, 2.6) -- (3.46, 2.6) -- cycle;
                \foreach \k in {0,...,8} \draw[fill=gray!15, draw=black!60] (0.15+\k*0.44, 0.2) rectangle ++(0.18, 2.4);
            \end{scope}
            \fill[naranja] (1.98, -0.22) circle (0.07);
            \node[font=\scriptsize] at (1.98, -0.5) {dock};
        \end{tikzpicture}
        \caption{Espina de pescado}
    \end{subfigure}
\end{figure}

\clearpage
\subsection{R3. Figura 2.4, las dos políticas de ruteo}
\donde{2.3.3 Order picking, misma posición.}
\begin{figure}
    \centering
    \begin{subfigure}{0.48\textwidth}
        \centering
        \begin{tikzpicture}[x=1cm, y=1cm]
            \foreach \i in {0,...,6} \draw[fill=gray!25, draw=black!60] (\i*0.62, 0.35) rectangle ++(0.24, 2.3);
            \draw[very thick, azul, -{Stealth}]
                (0.43, 0.15) -- (0.43, 2.85) -- (1.05, 2.85) -- (1.05, 0.15) -- (2.91, 0.15) -- (2.91, 2.85) -- (3.53, 2.85) -- (3.53, 0.15);
            \draw[thick, azul!60, dashed] (3.53, 0.15) -- (3.53, 0.02) -- (0.43, 0.02) -- (0.43, 0.15);
            \foreach \p in {(0.36, 1.9), (1.12, 0.9), (2.84, 1.4), (3.60, 2.1)} \fill[naranja] \p circle (0.07);
            \fill[black!70] (0.43, -0.2) circle (0.06);
            \node[font=\scriptsize] at (0.43, -0.45) {inicio};
        \end{tikzpicture}
        \caption{Política en forma de $S$}
    \end{subfigure}
    \hfill
    \begin{subfigure}{0.48\textwidth}
        \centering
        \begin{tikzpicture}[x=1cm, y=1cm]
            \foreach \i in {0,...,6} \draw[fill=gray!25, draw=black!60] (\i*0.62, 0.35) rectangle ++(0.24, 2.3);
            \draw[very thick, azul, -{Stealth}]
                (0.43, 0.15) -- (0.43, 1.9) -- (0.43, 0.15) -- (1.05, 0.15) -- (1.05, 0.9) -- (1.05, 0.15) -- (2.91, 0.15) -- (2.91, 1.4) -- (2.91, 0.15) -- (3.53, 0.15) -- (3.53, 2.1) -- (3.53, 0.15);
            \draw[thick, azul!60, dashed] (3.53, 0.15) -- (3.53, 0.02) -- (0.43, 0.02) -- (0.43, 0.15);
            \foreach \p in {(0.36, 1.9), (1.12, 0.9), (2.84, 1.4), (3.60, 2.1)} \fill[naranja] \p circle (0.07);
            \fill[black!70] (0.43, -0.2) circle (0.06);
            \node[font=\scriptsize] at (0.43, -0.45) {inicio};
        \end{tikzpicture}
        \caption{Política de retorno}
    \end{subfigure}
\end{figure}

\subsection{R4. Figura 2.6, el diente de sierra}
\donde{2.5.1 Dedicado o compartido, misma posición.}
\begin{figure}
    \centering
    \begin{tikzpicture}
        \begin{axis}[width=0.7\textwidth, height=4.2cm, xlabel={Tiempo}, ylabel={Stock en la ubicación}, xmin=0, xmax=3, ymin=0, ymax=1.15, xtick={0,1,2,3}, xticklabels={$0$,$T$,$2T$,$3T$}, ytick={0.5,1}, yticklabels={$Q/2$,$Q$}, tick label style={font=\scriptsize}, label style={font=\scriptsize}, axis lines=left]
            \addplot[very thick, azul] coordinates {(0,1) (1,0) (1,1) (2,0) (2,1) (3,0)};
            \addplot[dashed, naranja, thick, domain=0:3] {0.5};
        \end{axis}
    \end{tikzpicture}
\end{figure}

\subsection{R5. Figura 2.7, dos ubicaciones y su costo de acceso}
\donde{2.5.2 Las reglas de asignación, misma posición. Cada viaje es una sola
línea con punta en los dos extremos, ida y vuelta, y salen del \emph{dock} en
direcciones distintas para no pisarse. El pie sigue diciendo $c_A < c_B$.}
\begin{figure}
    \centering
    \begin{tikzpicture}[x=1cm, y=1cm]
        \foreach \i in {0,...,7} \draw[fill=gray!25, draw=black!60] (\i*0.62, 0.5) rectangle ++(0.24, 2.4);
        \fill[black!70] (2.29, 0.1) circle (0.07);
        \node[font=\scriptsize] at (2.29, -0.2) {dock};
        % A: viaje corto, vertical desde el dock
        \draw[very thick, azul, {Stealth}-{Stealth}] (2.29, 0.17) -- (2.29, 1.0);
        \fill[azul] (2.29, 1.0) circle (0.09);
        \node[font=\scriptsize, anchor=south] at (2.29, 1.12) {$A$};
        % B: viaje largo, sale del dock hacia la izquierda por el corredor frontal
        \draw[very thick, naranja, dashed, {Stealth}-{Stealth}] (2.2, 0.1) -- (0.43, 0.1) -- (0.43, 2.6);
        \fill[naranja] (0.43, 2.6) circle (0.09);
        \node[font=\scriptsize, anchor=south] at (0.43, 2.95) {$B$};
    \end{tikzpicture}
\end{figure}

\subsection{R6. Figura 2.5, el perfil esquemático}
\donde{2.4 Perfil de actividad y demanda, misma posición. Barras azul y
naranja, y la leyenda con un solo cuadrado por serie.}
\begin{figure}
    \centering
    \begin{subfigure}{0.48\textwidth}
        \centering
        \begin{tikzpicture}
            \begin{axis}[width=\linewidth, height=4.6cm, xlabel={Fracción de las SKU, ordenadas por picks}, ylabel={Fracción acumulada de picks}, xmin=0, xmax=1, ymin=0, ymax=1, xtick={0,0.2,0.5,1}, ytick={0,0.5,0.8,1}, tick label style={font=\scriptsize}, label style={font=\scriptsize}, grid=none, axis lines=left]
                \addplot[very thick, azul, domain=0:1, samples=100] {(1-exp(-8*x))/(1-exp(-8))};
                \draw[dashed, naranja] (axis cs:0.2,0) -- (axis cs:0.2,0.8) -- (axis cs:0,0.8);
            \end{axis}
        \end{tikzpicture}
        \caption{Curva de Pareto de los picks}
    \end{subfigure}
    \hfill
    \begin{subfigure}{0.48\textwidth}
        \centering
        \begin{tikzpicture}
            \begin{axis}[width=\linewidth, height=4.6cm, ybar, bar width=4pt, xlabel={Líneas por pedido}, ylabel={Fracción}, xmin=0.4, xmax=6.6, ymin=0, ymax=0.7, xtick={1,...,6}, ytick={0,0.2,0.4,0.6}, tick label style={font=\scriptsize}, label style={font=\scriptsize}, legend style={font=\scriptsize, draw=none, at={(0.97,0.95)}, anchor=north east}, legend image code/.code={\draw[#1] (0cm,-0.08cm) rectangle (0.25cm,0.08cm);}, axis lines=left]
                \addplot[fill=azul!45, draw=azul] coordinates {(1,0.60) (2,0.20) (3,0.10) (4,0.05) (5,0.03) (6,0.02)};
                \addplot[fill=naranja, draw=naranja!70!black] coordinates {(1,0.34) (2,0.23) (3,0.17) (4,0.11) (5,0.08) (6,0.07)};
                \legend{De los pedidos,De los picks}
            \end{axis}
        \end{tikzpicture}
        \caption{Distribución de líneas por pedido}
    \end{subfigure}
\end{figure}

% =====================================================================
\clearpage
\section{Descartadas y pendientes}
% =====================================================================
\begin{itemize}
    \item \textbf{Quién se mueve}, \textbf{mapa del capítulo} y el
    \textbf{ejemplo numérico del término cuadrático}, descartadas.
    \item \textbf{Perfil con datos}, va al capítulo de los datos. Queda anotado
    que en \texttt{data/v2/medium} las líneas por orden son uniformes entre 2 y
    5 por construcción del simulador.
    \item \textbf{Foto de un depósito real}, la aportás vos.
    \item \textbf{Tabla 2.1}, se rehace aparte.
\end{itemize}

\end{document}
